\documentclass[10pt,a4paper]{amsart}
\usepackage[margin=19mm]{geometry}
\usepackage{amsmath,amssymb,mathtools}
\usepackage[scr=rsfs,cal=euler]{mathalfa}
\usepackage{xfrac}
\usepackage[unicode,hidelinks,bookmarks=false]{hyperref}
\newcommand{\ie}{i.e.\ }
\newcommand{\cf}{cf.\ }
\newcommand{\resp}{respectively\ }
\newcommand{\Subsection}{\subsection}
\providecommand{\nobreakdash}{\nobreak\hbox{-}}
\setlength{\emergencystretch}{2em}
\multlinegap0pt
\usepackage{amssymb}
\usepackage{mathtools}
\usepackage[shortlabels]{enumitem}
\newtheorem{lemma}{Lemma}
\newtheorem{theorem}{Theorem}
\newcommand{\cP}{\mathcal P}
\DeclareMathOperator{\supp}{supp}
\title{Countable IET Models and Defect Sets for Interval Translation Maps}
\author{Sergey Kryzhevich, Khosro Tajbakhsh, Reza Yaghmaeian}
\date{Source: arXiv:2609.01670v1 (CC BY 4.0)}
\begin{document}
\maketitle

\noindent\emph{Source and licence.} Countable IET Models and Defect Sets for Interval Translation Maps. Version arXiv:2609.01670v1 (CC BY 4.0), distributed by arXiv under the Creative Commons Attribution 4.0 International licence, http://creativecommons.org/licenses/by/4.0/, as recorded in the arXiv licence metadata for this exact version and shown on the abstract page https://arxiv.org/abs/2609.01670v1. Full licence notice: \texttt{licenses/arxiv-2609.01670-CC-BY-4.0.txt}.\par\smallskip

\noindent\textit{Editorial scope.} Counted unit: the article's theorem thm:ae-invertible (source0.tex lines 790--796). The article's complete original proof of it is delivered with it (source0.tex lines 799--1007, one \texttt{proof} environment of 209 lines). No mathematics is written, repaired or re-derived: the statement and the proof are the article's own lines, in the article's own notation, and no retained line is substituted or re-flowed.  Retained uncounted as the source-local context the proof needs: lines 708--713; lines 758--768; lines 775--784.  Nothing was omitted from any retained span, so the package carries no omission record.  No retained line contains a citation, so the wrapper carries no bibliography and no bibliography entry.  No retained line contains a figure environment, an image-inclusion command or drawing code, so no image is delivered and the package declares no asset dependency.  Editorial formatting only, in addition: an AMS \texttt{amsart} preamble that restates the article's own declarations and macros used by the retained lines, namely the environments \texttt{lemma}, \texttt{theorem} on separate counters, together with the standard packages those lines need.  The retained environments print in this document as Lemma 1, Theorem 1 in order of appearance, whereas the article prints the same items with the numbering of its own full document.  The complete native source of the article as distributed in the arXiv e-print for this exact version is bundled unchanged as \texttt{sources/209134/source0.tex}.
\begin{lemma}[Zero entropy of the continuity partition]\label{lem:zero-partition-entropy}
Let $\mu$ be an $S$-invariant Borel probability measure. Then
\[
        h_\mu(S,\cP)=0.
\]
\end{lemma}

\begin{lemma}[The branch is a function of the image point]\label{lem:branch-image}
There exists a measurable function
\[
        \beta:I\to\{1,\ldots,n\}
\]
such that
\[
        \alpha(x)=\beta(Sx)
\]
for $\mu$-almost every $x$.
\end{lemma}

\begin{lemma}[Automorphic core from two one-sided inverses]\label{lem:automorphic-core}
Let \((Z,\mathcal A,\mu)\) be a standard Borel probability space and let
\(S,R:Z\to Z\) be measurable maps. Assume that \(S\) preserves \(\mu\), and that
\[
        R\circ S=\mathrm{id},\qquad S\circ R=\mathrm{id}
\]
hold \(\mu\)-almost everywhere. Then there exists a full-measure Borel set
\(Z_0\subset Z\) such that \(S(Z_0)=Z_0\), \(R(Z_0)=Z_0\), and the restrictions
\(S|_{Z_0}\) and \(R|_{Z_0}\) are inverse measurable measure-preserving bijections.
\end{lemma}

\begin{theorem}[Almost-everywhere invertibility]\label{thm:ae-invertible}
Let $S$ be an interval translation map and let $\mu$ be a non-atomic $S$-invariant Borel probability measure. Put $J=\supp\mu$. Then $S$ is invertible $\mu$-almost everywhere on $J$. More precisely, there exists an invariant full-measure Borel set $X_0\subset J$ such that
\[
        S:X_0\to X_0
\]
is a measurable bijection whose inverse is measurable and measure-preserving.
\end{theorem}

\begin{proof}
By Lemma~\ref{lem:zero-partition-entropy}, \(h_\mu(S,\cP)=0\). Hence
Lemma~\ref{lem:branch-image} gives a measurable function
\(\beta:I\to\{1,\ldots,n\}\) such that
\[
        \alpha(x)=\beta(Sx)
\]
for \(\mu\)-almost every \(x\).

Choose a fixed point \(x_*\in I\). For every branch \(I_k\), define a measurable map
\(R_k:I\to I\) by
\[
R_k(y)=
\begin{cases}
        y-c_k, & y\in S(I_k),\\
        x_*,   & y\notin S(I_k).
\end{cases}
\]
On \(S(I_k)\), this is the inverse of the restriction \(S|_{I_k}\). Define
\[
        R(y)=R_{\beta(y)}(y).
\]
Since \(\beta\) has finite range and each \(R_k\) is measurable, \(R:I\to I\) is
measurable.

Let \(E\subset J\) be a Borel full-measure set on which
\(\alpha(x)=\beta(Sx)\). If \(x\in E\cap I_k\), then \(\beta(Sx)=k\) and
\(Sx\in S(I_k)\), so
\[
        R(Sx)=R_k(Sx)=x.
\]
Thus \(R\circ S=\mathrm{id}\) \(\mu\)-almost everywhere on \(J\).

Next set
\[
        F=\{y\in I:S(Ry)=y\}.
\]
The set \(F\) is Borel. If \(x\in E\), then
\(S(R(Sx))=Sx\), so \(Sx\in F\), equivalently \(E\subset S^{-1}F\). By
invariance,
\[
        \mu(F)=\mu(S^{-1}F)\geq\mu(E)=1.
\]
Hence \(S\circ R=\mathrm{id}\) \(\mu\)-almost everywhere.

We now replace \(S\) and \(R\), on null sets only, by measurable
self-maps of \(J\). Put
\[
        N_S
        :=
        \{x\in J:Sx\notin J\}.
\]
Since \(\mu(J)=1\) and \(S\) preserves \(\mu\),
\[
        \mu(N_S)
        =
        \mu\bigl(J\setminus S^{-1}J\bigr)
        =
        0.
\]

Moreover,
\[
        S(E)
        =
        \bigcup_{k=1}^{n}\tau_k(E\cap I_k)
\]
is Borel. Since
\[
        E\subset S^{-1}(S(E)),
\]
invariance gives
\[
        \mu(S(E))
        =
        \mu\bigl(S^{-1}(S(E))\bigr)
        \geq
        \mu(E)
        =
        1.
\]
Thus \(S(E)\cap J\) has full measure in \(J\). For every
\(y=Sx\) with \(x\in E\), we have
\[
        R(y)=R(Sx)=x\in J.
\]
Consequently, the Borel set
\[
        N_R
        :=
        \{y\in J:R(y)\notin J\}
\]
is \(\mu\)-null.

Choose a fixed point \(z_*\in J\) and define
\[
        \widetilde S(x)
        :=
        \begin{cases}
        S(x),&x\in J\setminus N_S,\\
        z_*,&x\in N_S,
        \end{cases}
\]
and
\[
        \widetilde R(y)
        :=
        \begin{cases}
        R(y),&y\in J\setminus N_R,\\
        z_*,&y\in N_R.
        \end{cases}
\]
Then \(\widetilde S,\widetilde R:J\to J\) are measurable.

We verify explicitly that the two almost-everywhere inverse identities
are unchanged. Let
\[
        A
        :=
        \{x\in J:R(Sx)=x\}.
\]
The set \(A\) has full measure. If
\[
        x\notin
        (J\setminus A)\cup N_S\cup S^{-1}(N_R),
\]
then
\[
        \widetilde S(x)=S(x),
        \qquad
        \widetilde R(Sx)=R(Sx),
\]
and hence
\[
        \widetilde R(\widetilde Sx)=x.
\]
Now
\[
        \mu(S^{-1}(N_R))=\mu(N_R)=0
\]
by the \(S\)-invariance of \(\mu\). Therefore
\[
        \widetilde R\circ\widetilde S
        =
        \operatorname{id}
        \qquad
        \mu\text{-almost everywhere on }J.
\]

Similarly, let
\[
        B
        :=
        \{y\in J:S(Ry)=y\}.
\]
The set \(B\) has full measure. If
\[
        y\in B\setminus N_R,
\]
then \(R(y)\in J\) and
\[
        S(Ry)=y\in J.
\]
It follows that \(R(y)\notin N_S\). Hence neither modification affects
the composition at \(y\), and
\[
        \widetilde S(\widetilde Ry)
        =
        S(Ry)
        =
        y.
\]
Thus
\[
        \widetilde S\circ\widetilde R
        =
        \operatorname{id}
        \qquad
        \mu\text{-almost everywhere on }J.
\]

Finally, \(\widetilde S\) still preserves \(\mu\). Indeed, for every
Borel set \(C\subset J\),
\[
        \widetilde S^{-1}(C)
        \mathbin{\triangle}
        \bigl(S^{-1}(C)\cap J\bigr)
        \subset N_S,
\]
and therefore
\[
        \mu(\widetilde S^{-1}(C))
        =
        \mu(S^{-1}(C))
        =
        \mu(C).
\]
Renaming \(\widetilde S\) and \(\widetilde R\) as \(S\) and \(R\), we
may therefore apply Lemma~\ref{lem:automorphic-core} on the standard
Borel probability space \((J,\mu)\).

Lemma~\ref{lem:automorphic-core} gives an invariant full-measure Borel
set \(X_0\subset J\) on which \(S\) and \(R\) are inverse measurable
bijections. Since \(S\) preserves \(\mu\), the bijection
\[
        S:X_0\longrightarrow X_0
\]
is measure preserving.
\end{proof}

\end{document}
